\chapter{Элементарный объём и материя}
\label{ch:matter}

\section{Бинарная занятость и планкон}
\label{sec:planckon-floor0}

На каждой ячейке~$\PlanckCell$ вводится бинарная занятость $b(x)\in\{0,1\}$:
\textbf{планковская дырка} ($b=0$) или \textbf{планкон} ($b=1$).
Плотность материи принимает значения $0$ или планковская~$\rho_P$.
Первичный дефект --- одна ячейка.

\section{Несжимаемая решёточная жидкость}

\begin{definition}[Модуль упругости $K_P$]
\label{def:kp}
$K_P=u_P$ --- планковское давление жидкости;
вращатель~$\Phi$ насыщается при $\rho\to u_P$.
\end{definition}

Отклонение плотности от фона вызывает упругий фазовый отклик;
избыток энергии уходит волной ($n=0$) или вихрем.
Уравнения Эйлера и Навье--Стокса на~$T$ --- макроосреднение.

\section{Теорема о механической лестнице}
\label{sec:theorem-51}

\begin{theorem}[Лестница Ландау на ячейке]
\label{thm:51}
После \cref{thm:discreteness,thm:g-law} локальные механические величины на~$\hv$ образуют
единственную лестницу $s_0\to p_0\to L_0\to F_0\to E_0$
без трёх независимых подгоняемых параметров.
\end{theorem}

\begin{lemma}[5.1a]
\label{lem:51a}
\[
  s_0=\hbar/2.
\]
\end{lemma}

\begin{proof}
Минимальный фазовый шаг на такт
\[
  \Delta\phi_{\min}=\tfrac12.
\]
Ненулевой осциллирующий фон задаёт минимальное событие изменения аргумента на~$\hv$ за такт~$\htick$:
\[
  s_0=\hbar\,\Delta\phi_{\min}=\hbar/2.
\]
\end{proof}

\begin{lemma}[5.1b]
\label{lem:51b}
Из $s_0$ и пары $(\hl,\htick)$: $p_0=s_0/\hl$, $E_0=s_0/\htick$, $m_{\arg}=E_0/c^2$,
и потому $p_0=m_{\arg} c_0/2$.
\end{lemma}

\begin{proof}
Квант действия $s_0$, отнесённый к шагу, есть импульс; тот же квант, отнесённый к такту, есть энергия:
\[
  p_0=\frac{s_0}{\hl}=\frac{\hbar}{2\hl},\qquad
  E_0=\frac{s_0}{\htick}.
\]
Тактовая скорость $c_0=\hl/\htick$ связывает их: $E_0=p_0 c_0$.
Масса аргумента определяется через макроскопическую скорость света, уже выведенную из оболочки:
$c=\kgeom c_0$ при $\kgeom=1/\sqrt{2}$, то есть $c^2=c_0^2/2$.
\[
  m_{\arg}=\frac{E_0}{c^2}=\frac{E_0}{c_0^2/2}=\frac{2E_0}{c_0^2}=\frac{2p_0}{c_0}.
\]
Отсюда $p_0=m_{\arg} c_0/2$.
Множитель $1/2$ --- это $\kgeom^2$.
\end{proof}

\begin{lemma}[5.1c]
\label{lem:51c}
$L_0=p_0\hl=s_0$, $g_M=c_0/(2\htick)$, $F_0=p_0/\htick=m_{\arg} g_M$ ---
соотношение $F=mg$ на одной ячейке~$\hv$.
\end{lemma}

\begin{proof}
Момент на одном шаге совпадает с действием:
\[
  L_0=p_0\hl=s_0.
\]
Из \cref{lem:51b} скорость, стоящая в $p_0=m_{\arg} v$, равна $v=c_0/2$.
Набрать её за один такт значит иметь ускорение
\[
  g_M=\frac{c_0/2}{\htick}=\frac{c_0}{2\htick}.
\]
Сила, меняющая импульс на один квант за один такт,
\[
  F_0=\frac{p_0}{\htick}.
\]
Подставляем $p_0=m_{\arg} c_0/2$:
\[
  m_{\arg} g_M=m_{\arg}\cdot\frac{c_0}{2\htick}=\frac{p_0}{\htick}=F_0.
\]
Это и есть $F=mg$ на одной ячейке: ни масса, ни ускорение не вводятся отдельно от~$s_0$.
\end{proof}

\begin{lemma}[5.1d]
\label{lem:51d}
Ударный член $\lfloor\mathcal{N}\rfloor$ на~$\mathbb{Z}$ $\Rightarrow$ $\Delta p\in p_0\cdot\mathbb{Z}$;
интеграл $\int\Im(z^*\nabla z)$ --- только макроописание на~$T$.
\end{lemma}

\begin{proof}
В замкнутом~$g$ (\cref{thm:g-law}) импульс на связи дискретен: $p_0=s_0/\hl$.
Удар $\lfloor\mathcal{N}\rfloor$ --- целое число фазовых квантов $\Delta\phi_{\mathrm{disc}}$
на такт (\cref{sec:congruence}, CL-4), поэтому приращение импульса
$\Delta p=p_0\cdot n$ с $n\in\mathbb{Z}$.

\smallskip
Интеграл $\int\Im(z^*\nabla z)\,\mathrm{d}^3x$ предполагает гладкое поле и меру Лебега ---
на~$M$ есть только суммы по узлам и разности по~$N(x)$.
После осреднения~$\mathcal{B}$ (\cref{def:soft-readout}) дискретные потоки сходятся
к маделунговому току на~$T$; на~$M$ интеграл не определён.
\end{proof}

\begin{lemma}[5.1e]
\label{lem:51e}
$E_0=p_0 c_0=F_0\hl=L_0/\htick$, $E=n_E E_0$, $n_E\in\mathbb{Z}$.
\end{lemma}

\begin{proof}
Из \cref{lem:51b,lem:51c} и $c_0=\hl/\htick$:
\begin{align*}
  E_0&=\frac{s_0}{\htick}=\frac{p_0\hl}{\htick}=p_0 c_0,\\
  F_0\hl&=\frac{p_0}{\htick}\cdot\hl=p_0 c_0,\\
  \frac{L_0}{\htick}&=\frac{s_0}{\htick}=E_0.
\end{align*}
Три записи совпадают, поэтому лестница замыкается на одном кванте энергии.
Число событий за такт есть целое число фазовых квантов импульса,
$n_E=\lfloor|\Phi|/\Delta\phi_{\mathrm{disc}}\rfloor$ (CL-4, \cref{sec:congruence}),
и энергия события $E=n_E E_0$ кратна тому же~$E_0$.
\end{proof}

\begin{proof}[Доказательство теоремы]
Минимальное событие на ячейке несёт действие $s_0=\hbar/2$ (\cref{lem:51a}).
Деление на $\hl$ и на $\htick$ даёт импульс и энергию;
масса через $c=\kgeom c_0$ превращает их в $p_0=m_{\arg} c_0/2$ (\cref{lem:51b}).
Сила одного кванта импульса за такт совпадает с $m_{\arg} g_M$ (\cref{lem:51c}).
Приращение импульса на решётке целочисленно, а непрерывный интеграл $\int\Im(z^*\nabla z)$ на~$M$ не определён (\cref{lem:51d}).
Три записи энергии $p_0 c_0$, $F_0\hl$ и $L_0/\htick$ равны одному и тому же~$E_0$ (\cref{lem:51e}).
Свободного параметра, который можно было бы подогнать отдельно для импульса, силы и энергии, не остаётся.
\end{proof}

\begin{corollary}[Закон Ньютона на ячейке]
$p_0=m_{\arg}(c_0/2)$, $F_0=m_{\arg} g_M$, $\Delta F=\Delta m\cdot g$ при $\Delta m\in m_{\arg}\mathbb{Z}$.
\end{corollary}

\section{Термодинамические потенциалы на \texorpdfstring{$T$}{T}}
\label{sec:thermo-potentials}

Классические потенциалы относятся к макроуровню~$T$: на~$M$ нет температуры прибора
и нет гиббсовского ансамбля по~$g$ (\cref{prop:determinism}).
Их можно задать однозначно, если зафиксировать абсолютную энтропию по Планку
(\cref{ch:axiom-rationale}) и связать макроскопические $U$, $P$, $V$ с лестницей
\cref{thm:51} после осреднения~\cref{def:soft-readout}.

\paragraph{Внутренняя энергия.}
По \cref{lem:51e} энергия на~$M$ дискретна: $E=\sum_x n_E(x)\,E_0$,
$n_E\in\mathbb{Z}_{\ge 0}$.
На~$T$ в объёме~$V$ прибор видит гладкую плотность возбуждений
$\mathscr{n}_E(X)=\mathcal{B}*n_E$ и определяет
\begin{equation}
\label{eq:internal-energy}
  U(V)=\int_V \mathscr{n}_E(X)\,E_0\,\mathrm{d}^3X
  = \int_V u(X)\,\mathrm{d}^3X,
  \qquad u=\mathscr{n}_E E_0.
\end{equation}
Вакуум ($n_E=0$ везде) даёт $U=0$; это согласуется с $S\to 0$ при $T\to 0$,
если основное состояние невырождено ($W=1$ в формулировке Планка).

\paragraph{Давление и модуль.}
Модуль~$K_P=u_P$ (\cref{def:kp}) --- планковское давление несжимаемой решёточной жидкости:
вращатель насыщается при $\rho_E\to K_P$.
В термодинамике объём~$V$ --- макрообъём прибора; давление~$P$ --- сопряжённая
к~$V$ величина при медленных деформациях фона.
В натуральных единицах ($c_0=\hl/\htick=1$) равновесный вакуум ($\rho_E=0$)
лежит у нижней границы насыщения; для малых отклонений $\delta\rho=\rho_E\ll K_P$
\begin{equation}
\label{eq:pressure-linear}
  P \;\approx\; K_P\,\frac{\delta\rho}{K_P} = \delta\rho,
  \qquad
  \left.\frac{\partial P}{\partial V}\right|_{T}
  \sim -\frac{K_P}{V}\quad\text{при фиксированной массе.}
\end{equation}
Полная нелинейная связь $P(\rho_E)$ задаётся тем же saturating gate,
из которого в \cref{sec:gl-free-energy} получен член $\tfrac{\beta}{2}|\Phi|^4$
($\beta\equiv\omega_{\mathrm{macro}}$);
$\mathcal{F}_{\mathrm{GL}}$ описывает \emph{поле}~$\Phi$, а $P$ и~$U$ --- скаляры
субъекта наблюдения на~$T$.

\paragraph{Энтропия.}
На~$M$ энтропия фазы --- мера смешения аргументов~$z$
в пределах одной ячейки~$\hv$ при фиксированном~$Q$.
После~$\mathcal{B}$ прибор не различает микрофазу внутри вокселя:
макроэнтропия~$S$ --- функция $T$ и макропараметров, согласованная с
\[
  S = k_{\mathrm{B}}\ln W
\]
(Планк), где $W$ --- число равновесных \emph{макро}состояний, неразличимых readout'ом.
Условие $S\xrightarrow[T\to 0]{}0$ задаёт нулевую точку шкалы и делает
абсолютными потенциалы ниже.

\paragraph{Функции Гельмгольца и Гиббса.}
При фиксированных $T$ и~$V$ (канонический ансамбль на~$T$) определяют
\begin{equation}
\label{eq:helmholtz}
  F = U - TS.
\end{equation}
При фиксированных $T$ и~$P$ вводят энтальпию $H=U+PV$ и
\begin{equation}
\label{eq:gibbs}
  G = H - TS = U + PV - TS = F + PV.
\end{equation}
Связи $S=-(\partial F/\partial T)_V$, $P=-(\partial F/\partial V)_T$ --- стандартные
тождества при гладкости $U(T,V)$; модель не подменяет их новыми аксиомами,
а фиксирует \emph{микроскопические} якоря $E_0$ и~$K_P$, из которых $U$ и~$P$
следуют после~\eqref{eq:internal-energy} и насыщения gate.
Полевая свободная энергия~\eqref{eq:gl-free-energy} --- другой уровень описания:
$\mathcal{F}_{\mathrm{GL}}[\Phi]$ для неоднородного $\Phi(X)$, тогда как $F$ в~\eqref{eq:helmholtz}
--- скаляр равновесного объёма~$V$.
Химический потенциал $\mu=(\partial F/\partial N)_T$.

\section{Симметрии и дискретные заряды}

\begin{proposition}[Теорема Нётер на решётке]
На~$M$ сохраняются:
\begin{itemize}
\item заряд при $U(1)$: $Q=\sum|z|^2$;
\item импульс: $\mathrm{div}\,\pi=0$ по модулю~$p_0$;
\item момент при $SO(2)$ изотropного~$N$: $L_z\in L_0\cdot\mathbb{Z}$;
\item энергия: $E\in E_0\cdot\mathbb{Z}$;
\item топологический заряд: $n\in\mathbb{Z}$.
\end{itemize}
\end{proposition}

\section{Внутренний спектр планкона}
\label{sec:planckon-internal-spectrum}

На одном~$\hv$ (\cref{sec:planckon-floor0}, этаж~0) два внутренних регистра,
не отдельные поля Стандартной модели:
фазовое кольцо $\Phi\in\mathbb{Z}_{N_{\mathrm{ring}}}$ ($N_{\mathrm{ring}}=512$)
и спинор~$SU(2)$ на ядре.

\begin{proposition}[Лестница $n_E$ на кольце]
\label{prop:brick-ne-ladder}
Минимальный скачок $\Delta\varphi_{\min}$ соответствует одному кванту~$E_0$:
41~тик $\Rightarrow$ $n_E=1$; Pauli-kick $=256=N_{\mathrm{ring}}/2$;
полный оборот $512$ тик $\Rightarrow$ $n_E=12$.
\end{proposition}

\begin{proposition}[Основное состояние планкон]
\label{prop:brick-vortex-ground}
Устойчивый vortex с $n=\pm1$ после релаксации:
$b=1$, $|\Delta\varphi|_{\mathrm{core}}<\Delta\varphi_{\min}$,
$n_E=0$ в gate-ledger, устойчивый вектор Блоха на ядре.
Вращение вокруг оси Блоха: $2\pi\to-1$, $4\pi\to+1$ (спин~$\tfrac12$).
\end{proposition}

Возбуждения с $n_E\ge 1$ в свободной эволюции стандартных seeds
пока не закрыты sim'ом --- отдельный под-leaf этажа~0.

\section{Геометрический квант связи \texorpdfstring{$\kappa_{\mathrm{link}}=1/|N|$}{kappa=1/|N|}}

\begin{proposition}
На FCC $|N|=12$ $\Rightarrow$ $\kappa_{\mathrm{link}}=1/12=\gamma=\nu_{\mathrm{CA}}/(c_0\hl)$.
На гексагональном срезе $|N|=6$ $\Rightarrow$ $1/6$.
\end{proposition}

Спектральные методы калибруют макроуровень~$T$.

\section{Фононы и спектр}

Периодичность~$\Lambda$ $\Rightarrow$ первая зона Бриллюэна;
$p=n p_0+\hbar k$, $k$ дискретен.
Фонон вакуума:, $v_a=2c_0$.
\[
\omega(k)\approx v_a|k|
\]
Фонон вещества --- второй кристалл над микроуровнем (гл.~\ref{ch:carrier}).
